\section{Conclusion}

Sample covariance matrices can produce unstable minimum-variance portfolio weights. This essay investigated whether shrinkage provides a more reliable alternative by using shrinkage targets and intensities to reduce sampling error. The principle of shrinkage is to trade some bias for lower estimation variance. Linear shrinkage applies the same intensity to all eigenvalues, whereas nonlinear shrinkage applies different adjustments depending on their estimated distortion. To examine whether the findings of Ledoit and Wolf extend to the South African market, the empirical application compared the sample covariance estimator, linear shrinkage towards a scaled-identity target and nonlinear shrinkage using selected JSE-listed equities. Linear shrinkage produced the most stable covariance matrix, slightly less extreme weights and the strongest out-of-sample performance, although its improvements were marginal. These small differences are consistent with using only 17 assets relative to 1,768 estimation observations, giving an \(N/T\) ratio of approximately 0.0096. The results suggest that shrinkage can be a useful risk-management tool, particularly in higher-dimensional settings, but does not guarantee superior portfolio performance.
